\scanpage{110}\mspage{\handwritten{$9'$}}

infinitely many ope\deleted{\unsure{a}}rations standing for the single ope\unsure{rat}ion $v\overset{i}{\circ}u$ say, which could be viewed in effect in terms of an arbitrary $n$-sequence \handwritten{$(n\geq2)$} of ``composable'' $i$-objects $u_1=u,u_{\handwritten{2}}=v,u_{\handwritten{3}},\ldots\,u_n$. The natural way to meet this ``objection'' would be to restrict to \deleted{a} pairs $(f,g)$ which cannot be factored non trivially through another object $K'$ as
\diagram{\begin{tikzcd}[ampersand replacement=\&,column sep=large]
K \arrow[r] \& K' \arrow[r,shift left=1.4ex,"{f',g'}"] \arrow[r,shift right=1.4ex] \arrow[r,phantom,"{\deleted{F_i}}" description] \& F_i\; .
\end{tikzcd}}
But even with such restrictions, there remain a lot of redundancies -- and this again seems to me in the nature of things, namely that there is no really natural, ``most economic'' way for achieving condition (A), by a stepwise construction meeting condition (B). For instance, in $\underline{C}_1$ already we will have \deleted{the} not merely the compositions $v\overset{i}{\circ}u$, but at the same time simultaneous compositions
\[
u_n\handwritten{\circ}u_{n-1}\handwritten{\circ}\ldots\handwritten{\circ}u_1 \tag*{(*)}
\]
for ``composable'' sequences of $i$-objects $(i\handwritten{\geq}1)$, without reducing this (as is c\unsure{o|u|a}stomary) to iteration of the binary composition $v\circ u$. Of course using the binary composition, and more generally iteration of $n'$-ary compositions with $n'\handwritten{<}n$ (when $n\handwritten{\geq}3$), we get an impressive bunch of ope\deleted{\unsure{a}}rations in the '$n$ variables $u_1,\ldots,u_n$, serving the same purpose as $(*)$. All these will be tied up by homotopies in the next step $\underline{C}_2$. We would like to think of this set of homotopies in $\underline{C}_2$ as a kind of ``transitive system of isomorphisms'' (of associativity), now the transitivity relations one is looking for will be replaced by homotopies again between compositions of homotopies, which will enter in the picture with $\underline{C}_{\handwritten{3}}$, \inserted{\handwritten{etc.}} Here the infinite story is exemplified by the more familiar situation of the two ways in which one could define a``$\unsure{\otimes|x}$-composition with associativity'' in a category, \deleted{using either} starting either in terms of a binary operation, or \inserted{with} a bunch of $n$-ary operations -- with\handwritten{,} in both ca\handwritten{\unsure{ses}}, the associativity isomorphisms being an essential part of the structure. Here again, while it is generally (and quite validly) felt that the two points of view are equivalent\unsure{;|,} and both have their advantages and their drawbacks, still it is not true \inserted{I believe} that the \inserted{two} categories of algebraic structures ``category with associative $\unsure{\otimes|x}$-operation'', using one or the other definition, are equivalent. Here the story though of the relation between the two notions is a finite one, due to the fact that it is related to the notion of $2$-categories or $2$-groupoids, instead of $\infty$-groupoids as before \ldots

Thus I dont feel really like spending much energy in cutting down redundanci\unsure{es}, but prefer working with a notion of $\infty$-groupoid which remain\unsure{s} partly indeterminate, the main features being embodied in the conditions (A) and (B) and in the description of $\underline{C}_{\unsure{o|0}}$, without other specification.

% archive sheet number: 109
